\section{Simulation Setup}

The current paper scaffold is instantiated with the Rabot outdoor scenario defined in \texttt{scenarios/rabot.yaml}. The scene is derived from an \gls{osm} boundary around the KU Leuven Gent Campus Rabot environment. Candidate movable-\gls{ap} positions are generated along building walls at low height, while user mobility is simulated over the walkable outdoor space with pedestrian- and bicycle-like speed profiles.

The simulation pipeline used throughout the manuscript follows the same sequence as the repository workflow: scene construction, candidate generation, user-trajectory generation, path-based \gls{csi} extraction, strategy comparison, and postprocessing into publication-oriented figures. This section should remain descriptive and reproducible; interpretation belongs in the results and discussion sections.

\begin{table}[t]
    \centering
    \caption{Key Rabot scenario parameters used by the current manuscript scaffold.}
    \label{tab:rabot_setup}
    \begin{tabular}{ll}
        \toprule
        Parameter & Value \\
        \midrule
        Carrier frequency & \SI{3.5}{\giga\hertz} \\
        Bandwidth & \SI{100}{\mega\hertz} \\
        User height & \SI{1.5}{\meter} \\
        Number of users & 10 \\
        Simulation horizon & \SI{600}{\second} \\
        Time step & \SI{10}{\second} \\
        Movable \glspl{ap} & 6 \\
        Fixed \glspl{ap} & 0 \\
        Candidate wall height & \SI{1.5}{\meter} \\
        Candidate wall spacing & \SI{18}{\meter} \\
        Minimum candidate spacing & \SI{12}{\meter} \\
        Relocation window & \SI{60}{\second} \\
        Compared strategies & baseline, local heuristic \\
        Exact search & disabled in current Rabot config \\
        Coverage maps & disabled in current Rabot config \\
        \bottomrule
    \end{tabular}
\end{table}

The manuscript should also state the current output contract because the figures in \cref{sec:results_geometry,sec:results_sinr,sec:results_relocation} are directly linked to committed simulation artifacts. In the present workflow, the paper expects scene-layout figures in \texttt{outputs/rabot/} and postprocessed \gls{sinr}/schedule figures in \texttt{outputs/rabot/postprocessing/}.

\gilles{Optional later addition: add one compact pipeline schematic if the venue benefits from a visual overview.}
